%%  Thesis Progress Document
%%  Author : Ali Tayyar
%%  Status : Living document — updated incrementally
\documentclass[12pt, a4paper, twoside]{report}

%% ---------- Packages ----------
\usepackage[T1]{fontenc}
\usepackage[utf8]{inputenc}
\usepackage[a4paper, margin=2.5cm]{geometry}
\usepackage{amsmath, amssymb, amsthm}
\usepackage{mathtools}
\usepackage{bm}           % bold math
\usepackage{graphicx}
\usepackage{subcaption}
\usepackage{booktabs}
\usepackage{array}
\usepackage{multirow}
\usepackage{xcolor}
\usepackage{algorithm}
\usepackage{algpseudocode}
\usepackage{listings}
% hyperref MUST come before cleveref
\usepackage[colorlinks=true, linkcolor=blue, citecolor=blue, urlcolor=blue]{hyperref}
\usepackage{cleveref}
\usepackage{fancyhdr}
\usepackage{emptypage}
\usepackage{parskip}

%% ---------- Colours ----------
\definecolor{codeblue}{rgb}{0.13,0.29,0.53}
\definecolor{codegray}{rgb}{0.5,0.5,0.5}
\definecolor{codegreen}{rgb}{0.0,0.5,0.0}
\definecolor{backcolour}{rgb}{0.97,0.97,0.97}

%% ---------- Code listing style ----------
\lstdefinestyle{mystyle}{
    backgroundcolor=\color{backcolour},
    commentstyle=\color{codegreen},
    keywordstyle=\color{codeblue},
    numberstyle=\tiny\color{codegray},
    stringstyle=\color{orange},
    basicstyle=\ttfamily\footnotesize,
    breakatwhitespace=false,
    breaklines=true,
    captionpos=b,
    keepspaces=true,
    numbers=left,
    numbersep=5pt,
    showspaces=false,
    showstringspaces=false,
    showtabs=false,
    tabsize=2
}
\lstset{style=mystyle}

%% ---------- Theorem environments ----------
\theoremstyle{plain}
\newtheorem{theorem}{Theorem}[chapter]
\newtheorem{proposition}[theorem]{Proposition}
\newtheorem{lemma}[theorem]{Lemma}
\newtheorem{corollary}[theorem]{Corollary}

\theoremstyle{definition}
\newtheorem{definition}[theorem]{Definition}
\newtheorem{remark}[theorem]{Remark}

%% ---------- Hyperref setup ----------
\hypersetup{
    colorlinks=true,
    linkcolor=codeblue,
    citecolor=codegreen,
    urlcolor=codeblue,
}

%% ---------- Header / footer ----------
\pagestyle{fancy}
\fancyhf{}
\fancyhead[LE,RO]{\thepage}
\fancyhead[RE]{\textit{\nouppercase{\leftmark}}}
\fancyhead[LO]{\textit{\nouppercase{\rightmark}}}
\renewcommand{\headrulewidth}{0.4pt}

%% ---------- Notation macros ----------
\newcommand{\ux}{u^*}           % target value
\newcommand{\uhat}{\hat{u}^*}   % predicted value
\newcommand{\dxi}{\Delta x_i}
\newcommand{\dti}{\Delta t_i}
\newcommand{\bw}{\mathbf{w}}
\newcommand{\bu}{\mathbf{u}}
\newcommand{\bA}{\mathbf{A}}
\newcommand{\bb}{\mathbf{b}}
\newcommand{\bI}{\mathbf{I}}
\newcommand{\R}{\mathbb{R}}
\newcommand{\Lop}{\mathcal{L}}
\newcommand{\normone}[1]{\lVert #1 \rVert_1}
\newcommand{\normtwo}[1]{\lVert #1 \rVert_2}

%% ============================================================
\begin{document}
  

%% ---------- Title page ----------
\begin{titlepage}
    \centering
    \vspace*{2cm}
    {\LARGE\bfseries Reinforcement Learning for Adaptive
    PDE Solving:\\[0.4em]
    A Space--Time Mesh-Free Framework\par}
    \vspace{1.5cm}
    {\large Ali Tayyar\par}
    \vspace{0.5cm}
    {\large\itshape Thesis Progress Document\par}
    \vspace{0.3cm}
    {\large Version 1: \today\par}
    \vfill
    {\large Supervised by: Dr. Joseph Bakarji\par}
    \vspace{0.5cm}
    {\large American University of Beirut\par}
\end{titlepage}

%% ---------- Abstract ----------
\chapter*{Abstract}
\addcontentsline{toc}{chapter}{Abstract}

This document records the ongoing development of a framework for
solving partial differential equations (PDEs) using Reinforcement
Learning (RL) in the context of mesh-free local reconstruction.
The idea is to treat point placement in the space time
domain as a sequential decision problem: rather than predefining a
fixed grid, an agent learns \emph{where} to evaluate the solution
next starting from the initial condition (IC), guided by a reward signal derived from accuracy and speed of the integration.

The work reported in this document covers the first part of the project, which focuses on building the reconstruction \emph{environment} within which RL will operate within.
This environment is based on a Taylor-constrained local linear
predictor closely related to the Generalised Finite Difference
Method (GFDM).
We establish an error bound for the predictor, demonstrate that the method recovers classical finite difference (FD) stencils when used on a fixed grid, and validate a full space-time rollout on the advection-diffusion equations against the analytical and FD reference solutions.
Open questions regarding neighbour selection, computational cost,
and the RL formulation are identified.

\newpage

%% ---------- Table of contents ----------
\tableofcontents
\newpage
\listoffigures
\newpage
\chapter{Introduction}
\label{chap:intro}
\label{sec:motivation}
\section{Motivation}

Numerical simulation of partial differential equations (PDEs) is a cornerstone of scientific computing, with applications ranging from fluid and solid mechanics to transport phenomena. Classical Eulerian approaches such as finite difference, finite element, and finite volume methods rely on meshing our space-time domain using structured or unstructured meshes. 
While these methods are well-established and highly accurate, their performance critically depends on the chosen discretisation and mesh structure, often requiring fine grids and small time steps to resolve complex dynamics, which leads to substantial computational cost.

To mitigate this cost, adaptive strategies such as Adaptive Mesh Refinement (AMR) dynamically allocate resolution to regions of interest\cite{plewa2005}. In recent years, AMR has been casted as a sequential decision-making problem with the use of Reinforcement Learning, where the system will start with a coarse grid and start learning where to derefine or refine each cell, demonstrating competitive performance compared to classical heuristics \cite{yang2021rl_amr, FOUCART23}. Despite these advances, such approaches still operate within a predefined mesh structure and do not address the more fundamental question of whether the system can learn the mesh on its own.

Mesh-free methods provide a more flexible alternative by removing the need for a predefined grid \cite{griebel05}. Rather than discretising the domain with a mesh, these methods operate directly on scattered point clouds, approximating differential operators and function values through local stencils from the chosen neighbours. The Generalised Finite Difference Method (GFDM) constructs derivative approximations at each node by enforcing Taylor consistency constraints over a local neighbourhood \cite{benito2001, gavete2003}, and has been successfully applied to a wide range of PDE problems \cite{zheng22}. For function value reconstruction specifically, the same Taylor-moment framework underlies Moving Least Squares (MLS) approximation \cite{lancaster81}, which fits a local polynomial to neighbouring values and is the foundation of several meshless solvers \cite{gambaruto2026}. RBF-based methods offer an alternative reconstruction strategy with favourable approximation properties on scattered data \cite{fornberg2015}, though they introduce a shape parameter that requires careful tuning. These methods naturally support local adaptivity and handle irregular geometries without remeshing. However, in practice, the placement of evaluation points is still determined by heuristic or user-defined sampling strategies, and reconstruction accuracy is sensitive to the local point distribution.

The present work adopts this perspective and proposes a framework for solving PDEs in which point placement in space--time is learned via reinforcement learning. The approach is built on a mesh-free reconstruction environment based on Taylor-consistent local approximations with PDE constraints, closely related to GFDM. This enables a fundamentally different paradigm in which the solver jointly learns both the solution and the discretisation strategy, opening the door to more efficient and adaptive numerical methods.
Thus, the project can be divided into two tightly coupled components:

\begin{enumerate}
    \item \textbf{The reconstruction environment:}
    Given a set of previously evaluated space--time points and their solution values, compute the solution at a new target point using a local weighted stencil. Where it must be
    accurate, stable, and fast enough to serve as the step function inside an RL training loop.

    \item \textbf{The RL agent:} Given the current cloud of evaluated points and the
    reconstruction environment, learn a policy that selects the next point to evaluate so as to maximise solution accuracy while minimising the total number of evaluations.
    This component has not yet been implemented; the work
    reported here is entirely devoted to component~1.
\end{enumerate}
% ============================================================
\chapter{Mathematical Background}
\label{chap:theory}
%% ============================================================

\section{Problem Setup: Local Linear Prediction}
\label{sec:setup}

Let $u : \R^2 \to \R$ be a smooth scalar field defined on the
space--time domain $(x,t)$.
Given a target point $\hat{u}^* (x^*, t^*)$ and $n$ neighbouring
points $z_i = (x_i, t_i)$ with known values $u_i = u(z_i)$, we
seek a scalar approximation
\begin{equation}
    \uhat = \sum_{i=1}^{n} w_i\, u_i = \bw^\top \bu, 
    \label{eq:predictor}
\end{equation}
i.e. a linear transformation $\R^n \to \R$ that will take $n$ neighbour point values and return the solution at the target coordinate $z^*$, and the weights $w_i$ depend only on the geometry of the
neighbourhood, i.e.\ on the offsets
\begin{equation}
    \dxi = x_i - x^*, \qquad \dti = t_i - t^*.
\end{equation}
The prediction error is $e = \uhat - \ux$.

\section{Taylor Consistency Constraints}

Expanding $u_i$ in a Taylor series about $z^*$ up to order $p$
and substituting into \eqref{eq:predictor}, the error $e$ vanishes
Expand $u(z_i)$ around $z^*$:
\begin{equation}
u(z_i)=u^* + u_x^*\,\Delta x_i + u_t^*\,\Delta t_i
+\frac{1}{2}u_{xx}^*\,\Delta x_i^2 + u_{xt}^*\,\Delta x_i\Delta t_i
+\frac{1}{2}u_{tt}^*\,\Delta t_i^2 + \cdots+R_i
\end{equation}

Plugging into \eqref{eq:predictor} gives
\begin{equation*}
    \hat u(z^*) =
u(z^*)\sum_i w_i
+ u_x(z^*)\sum_i w_i\Delta x_i
+ u_t(z^*)\sum_i w_i\Delta t_i
+ \frac{1}{2}u_{xx}(z^*)\sum_i w_i\Delta x_i^2 + \cdots+\sum_i w_i R_i \label{eq: taylor_exp}
\end{equation*}

So if we impose the \emph{moment conditions}
\begin{equation}\label{eq:moments-value}
\sum_{i=1}^n w_i = 1,\qquad
\sum_{i=1}^n w_i\Delta x_i = 0,\qquad
\sum_{i=1}^n w_i\Delta t_i = 0,\qquad
\ldots
\end{equation}
then all the undesired Taylor terms cancel and $\hat u(z^*)$ becomes accurate.
\label{sec:taylor}

if and only if the weights satisfy the \emph{moment conditions}:
\begin{equation}
    \sum_{i=1}^n w_i\, \dxi^k\, \dti^{\,\ell} =
    \begin{cases} 1 & k = \ell = 0, \\ 0 & 1 \le k+\ell \le p.
    \end{cases}
    \label{eq:moments}
\end{equation}
These conditions ensure that the predictor exactly reproduces all
polynomials of degree $\le p$.

\subsection{Matrix formulation,}
\label{sec:matrix}

Let $m$ and $n$ be the number of Taylor terms and the number of neighbours, respectively:

Construct matrix $A\in\mathbb{R}^{m\times n}$ where:
\begin{itemize}
\item each column corresponds to a neighbour,
\item each row corresponds to a Taylor term.
\end{itemize}


\[
\underbrace{
\begin{bmatrix}
1 & 1 & \cdots & 1 \\
\Delta x_1 & \Delta x_2 & \cdots & \Delta x_n \\
\Delta t_1 & \Delta t_2 & \cdots & \Delta t_n \\
\Delta x_1^2 & \Delta x_2^2 & \cdots & \Delta x_n^2 \\
\Delta x_1\Delta t_1 & \Delta x_2\Delta t_2 & \cdots & \Delta x_n\Delta t_n \\
\Delta t_1^2 & \Delta t_2^2 & \cdots & \Delta t_n^2 \\
\vdots & \vdots & \ddots & \vdots
\end{bmatrix}
}_{A\;\text{(Taylor matrix)}}
\;
\underbrace{
\begin{bmatrix}
w_1\\
w_2\\
w_3\\
w_4\\
w_5\\
\vdots\\
w_n
\end{bmatrix}
}_{w\;\text{(weights)}}
=
\underbrace{
\begin{bmatrix}
1\\
0\\
0\\
0\\
0\\
0\\
\vdots
\end{bmatrix}
}_{b\;\text{(Taylor constraints)}} \label{taylor system}
\]

\subsection*{Case 1: $n>m$, i.e. more neighbours than terms}
 Mostly common in the GFDM method. Underdetermined system with infinitely many solutions. Most stable regime.
\subsection*{Case 2: $n=m$}
Determined system.
\[
Aw=b.
\]
Unique interpolation solution if invertible.  
Sensitive to geometry, if A is not full rank, not invertible, it'll be pushed toward the overdetermined regime.
\subsection*{Case 3: $n<m$}
Overconstrained system. Typically unstable for prediction.
\subsection*{SVD/pseudoinverse unifies all cases}
Compute the singular value decomposition (SVD):
\[
A = U\Sigma V^\top.
\]
The pseudoinverse is
\[
A^+ = V\Sigma^+U^\top,
\]
where $\Sigma^+$ inverts nonzero singular values.

Then the numerical solution used in your code is:
\begin{equation}\label{eq:pinv-solve}
\bm w = A^+ \bm b.
\end{equation}
This automatically gives:
\begin{itemize} 
\item determined: the exact solution (same as $A^{-1}b$ when invertible),
\item underdetermined: the \textbf{minimum-norm} exact solution,
\item overdetermined: the \textbf{least-squares} solution minimizing $\|Aw-b\|_2$.
\end{itemize}
This is why pseudoinverse is the safest baseline: it never forces you to switch formulas by hand.


\subsection{Unifying value and derivative reconstruction}
\label{sec:unify}

The same matrix $\bA$ governs the approximation of any linear
functional $\ell$ of $u$ at $z^*$:
\begin{equation}
    \ell(u)\big|_{z^*} \approx \sum_{i=1}^n w_i^{(\ell)}\, u_i,
    \qquad
    \bA \bw^{(\ell)} = \bb^{(\ell)},
    \quad
    b_k^{(\ell)} = (\ell\, p_k)(0,0).
\end{equation}
Value reconstruction corresponds to $\ell = \mathrm{id}$ (giving
$\bb = (1,0,\ldots)^\top$); spatial or temporal derivatives
correspond to different right-hand sides with the same $\bA$.
This is the fundamental insight connecting our predictor to GFDM.


Examples for the same basis:
\begin{itemize} 
\item For $L=\partial_x$ (spatial derivative):
\[
\bm b^{(x)}
= (0,1,0,0,0,0)^\top.
\]
\item For $L=\partial_t$:
\[
\bm b^{(t)} = (0,0,1,0,0,0)^\top.
\]
\item For $L=\partial_{xx}$:
\[
\bm b^{(xx)}
= (0,0,0,2,0,0)^\top,
\]
because $\partial_{xx}(\Delta x^2)=2$ and the others vanish.
\end{itemize}
\textbf{Same $A$}, different $\bm b$ $\Rightarrow$ different weights.

\section{Error Bound}

\label{sec:error_bound}
From eq.(\ref{eq: taylor_exp}), let: \[e \;=\; u^* - \hat{u}^* \;=\; -\sum_{i=1}^n w_i R_i\] 

We now derive a bound on the local reconstruction error. Under the Taylor 
constraints \eqref{eq:moments}, the zeroth- and first-order terms cancel 
exactly, leaving the error determined entirely by the second-order remainder:
\[
    e = u^* - \hat{u}^* = -\sum_{i=1}^n w_i R_i,
\]
where $R_i$ collects all terms of order two and higher in the Taylor expansion 
of $u_i$ about $z^*$. Expanding explicitly:
\[
    R_i = \frac{1}{2}u_{xx}\,\Delta x_i^2 
        + u_{xt}\,\Delta x_i\Delta t_i 
        + \frac{1}{2}u_{tt}\,\Delta t_i^2 
        + O(h^3).
\]
Assuming all points lie within a ball of radius $h$ around $z^*$, each 
offset satisfies $|\Delta x_i|,\,|\Delta t_i| \leq h$, so every second-order 
term is bounded by $Ch^2$ where 
$C = \max(|u_{xx}|, |u_{xt}|, |u_{tt}|)$ evaluated at $(x^*, t^*)$. 
Therefore $|R_i| \leq Ch^2$ for all $i$, and applying the triangle inequality:
\[
    |e| = \Bigl|\sum_i w_i R_i\Bigr| 
        \leq \sum_i |w_i|\,|R_i| 
        \leq \sum_i |w_i| \cdot Ch^2 
        = \|w\|_1 \cdot C \cdot h^2.
\]
This gives the local reconstruction error bound
\[
    \boxed{|u^* - \hat{u}^*| \leq \|w\|_1 \cdot C \cdot h^2,} \label{eq: error_bound}
\]
where $\|w\|_1 = \sum_i |w_i|$. The bound has a natural interpretation: 
accuracy depends on two independent factors. The geometric factor $Ch^2$ 
improves as neighbours are placed closer to the target, while the weight 
factor $\|w\|_1$ reflects the stability of the stencil, i.e., if the system produces large weights, it amplifies the reminder even if $h$ is relatively small.s


\begin{remark}
The bound generalises directly to any truncation order $p$. Keeping terms up
to order $p$ gives $m = \tfrac{(p+1)(p+2)}{2}$ constraints in 1D.
The first discarded terms are of order $p+1$, so
\[
  |u^* - \hat{u}^*| \;\leq\; \|w\|_1 \cdot C_p \cdot h^{p+1},
\]
where $C_p$ bounds the $(p+1)$-th order derivatives of $u$ near $(x^*,t^*)$.
The system remains underdetermined as long as $n > m$, i.e.\
$n > \tfrac{(p+1)(p+2)}{2}$.
\end{remark}

\begin{remark}
Adding the PDE as an extra constraint (see \cref{sec:pde_constraint})
shrinks the feasible set for $\bw$, which in general
\emph{increases} $\normone{\bw}$ (by the minimum-norm property of
the pseudoinverse), making the bound looser even though the stencil
is more physically informed.
\end{remark}
\begin{remark}
    For the numerical results for the error bound please check section (\ref{sec:exp2})
\end{remark}

\section{Incorporating the PDE as a Constraint}
\label{sec:pde_constraint}
The Taylor consistency constraints derived in Section~2.2 are PDE-agnostic: they
guarantee that the predictor reproduces polynomials up to a given degree, but they
treat all partial derivatives of $u$ as independent unknowns, and regardless what the PDE is, it will treat them all the same.
When $u$ satisfies a known PDE, however, the partial derivatives are \emph{not}
independent, and the PDE links them together.
Exploiting this linkage allows us to reduce the number of constraints needed and enforcing the weights to be PDE constrained.

We derive the constraint matrices explicitly for two cases: pure diffusion and
advection diffusion PDE.

% -------------------------------------------------------------------
\subsubsection*{Case 1: Pure Diffusion \texorpdfstring{($u_t = \alpha u_{xx}$)}{}}
% -------------------------------------------------------------------

Starting from the Taylor expansion of a neighbour value $u_i$ about the target
$z^* = (x^*, t^*)$:
\begin{equation}
    u_i = u^* + u_x \Delta x_i + u_t \Delta t_i
          + \tfrac{1}{2} u_{xx} \Delta x_i^2 + u_{xt} \Delta x_i \Delta t_i
          + \tfrac{1}{2} u_{tt} \Delta t_i^2 + \cdots
\end{equation}
Substituting into the predictor $\hat{u}^* = \sum_i w_i u_i$ and grouping
by derivative gives:
\begin{equation}
    \hat{u}^* = u^* \sum_i w_i
    + u_x    \sum_i w_i \Delta x_i
    + u_t    \sum_i w_i \Delta t_i
    + \tfrac{1}{2} u_{xx} \sum_i w_i \Delta x_i^2
    + \cdots
\end{equation}
For this to equal $u^*$ the coefficient of every independent derivative must
vanish.
Without any PDE knowledge there are four independent first- and second-order
derivatives ($u_x$, $u_t$, $u_{xx}$, $u_{xt}$, $u_{tt}$), requiring five constraints.

Now impose $u_t = \alpha u_{xx}$.
Substituting into the expansion:
\begin{align}
    \hat{u}^* &= u^* \sum_i w_i
    + u_x    \sum_i w_i \Delta x_i
    + \alpha u_{xx} \sum_i w_i \Delta t_i
    + \tfrac{1}{2} u_{xx} \sum_i w_i \Delta x_i^2
    + \cdots \\
    &= u^* \sum_i w_i
    + u_x \sum_i w_i \Delta x_i
    + u_{xx} \sum_i w_i \!\left(\alpha \Delta t_i + \tfrac{1}{2}\Delta x_i^2\right)
    + \cdots
\end{align}
Thus, our constraints become:e
\begin{align}
    \sum w_i&=1 \label{eq:diff_row1} \\
    \sum_i w_i \Delta x_i &= 0 \label{eq:diff_row2}\\
    \sum_i w_i \!\left(\alpha \Delta t_i + \tfrac{1}{2}\Delta x_i^2\right) &= 0
    \label{eq:diff_row3}
\end{align}
So, the full matrix becomes $R^{3 \times n}$ constraint matrix and right-hand side are:
\begin{equation}
    A = \begin{pmatrix}
        1 & 1 & \cdots & 1 \\
        \Delta x_1 & \Delta x_2 & \cdots & \Delta x_n \\
        \alpha \Delta t_1 + \tfrac{1}{2}\Delta x_1^2
        & \alpha \Delta t_2 + \tfrac{1}{2}\Delta x_2^2
        & \cdots
        & \alpha \Delta t_n + \tfrac{1}{2}\Delta x_n^2
    \end{pmatrix},
    \qquad
    b = \begin{pmatrix} 1 \\ 0 \\ 0 \end{pmatrix}.
\end{equation} 
\begin{remark}
    The bound is $O(h^2)$ because we have 3 rows total, which is sufficient
    to cancel all first-order terms after the PDE substitution.
    Achieving $O(h^4)$ would require enough rows to cancel all third-order
    terms as well, which demands at least 10 constraints in 2D space--time
    regardless of how the PDE is used.
\end{remark}

% -------------------------------------------------------------------
\subsubsection*{Case 2: Advection--Diffusion \texorpdfstring{($u_t = \alpha u_{xx} - c\, u_x$)}{}}
% -------------------------------------------------------------------

For the advection--diffusion equation the procedure is identical but the PDE
couples $u_t$ to \emph{both} $u_{xx}$ and $u_x$.
Substituting $u_t = \alpha u_{xx} - c\, u_x$ into the Taylor expansion and
grouping by $u_x$ and $u_{xx}$:

\begin{align}
    \hat{u}^* &= u^* \sum_i w_i
    + u_x \sum_i w_i \Delta x_i
    + (\alpha u_{xx} - c\, u_x) \sum_i w_i \Delta t_i
    + \tfrac{1}{2} u_{xx} \sum_i w_i \Delta x_i^2
    + \cdots \\
    &= u^* \sum_i w_i
    + u_x    \sum_i w_i \!\left(\Delta x_i - c\, \Delta t_i\right)
    + u_{xx} \sum_i w_i \!\left(\alpha \Delta t_i + \tfrac{1}{2}\Delta x_i^2\right)
    + \cdots
\end{align}

Again two independent derivatives remain.
Setting their coefficients to zero:
\begin{align}
    \sum_i w_i \!\left(\Delta x_i - c\, \Delta t_i\right) &= 0
    \label{eq:adv_row2}\\
    \sum_i w_i \!\left(\alpha \Delta t_i + \tfrac{1}{2}\Delta x_i^2\right) &= 0
    \label{eq:adv_row3}
\end{align}
The full constraint matrix becomes:
\begin{equation}
    A = \begin{pmatrix}
        1 & 1 & \cdots & 1 \\
        \Delta x_1 - c\,\Delta t_1
        & \Delta x_2 - c\,\Delta t_2
        & \cdots
        & \Delta x_n - c\,\Delta t_n \\
        \alpha \Delta t_1 + \tfrac{1}{2}\Delta x_1^2
        & \alpha \Delta t_2 + \tfrac{1}{2}\Delta x_2^2
        & \cdots
        & \alpha \Delta t_n + \tfrac{1}{2}\Delta x_n^2
    \end{pmatrix},
    \qquad
    b = \begin{pmatrix} 1 \\ 0 \\ 0 \end{pmatrix}.
\end{equation}
This is exactly the matrix used in the implementation (see Section~3.3.3).

\paragraph{Comparison with the augmented-row formulation.}
Thus, we can think of the matrix $\tilde{A}$ as a combination of the taylor terms in addition to the PDE constraint:
\begin{equation}
    \tilde{A} = \begin{pmatrix} A_{\text{Taylor}} \\ \text{PDE row} \end{pmatrix}
    \in \mathbb{R}^{(m+1)\times n}.
\end{equation} 

\begin{remark}
    The weight factor $\|w\|_1$ is the dominant source of practical error
    growth: even though $O(h^2)$ convergence is guaranteed in order,
    a poorly conditioned stencil geometry can make $\|w\|_1$ large and
    effectively destroy the bound.
    This is discussed later in the computational sections.
\end{remark}
 
\subsubsection*{Exercise: Recovering the Explicit Heat Equation Scheme 
via Constrained Reconstruction} 
Given the heat equation $u_t = \alpha u_{xx}$, we reconstruct $u(x_*, t_*)$ from three values at time $t_* - \Delta t$:
\[
u_1 = u(x_* - \Delta x, t_* - \Delta t), \quad
u_2 = u(x_*, t_* - \Delta t), \quad
u_3 = u(x_* + \Delta x, t_* - \Delta t).
\]
We seek weights $(w_1, w_2, w_3)$ such that:
\[
u(x_*, t_*) \approx w_1 u_1 + w_2 u_2 + w_3 u_3.
\]

For this to equal $u_*$, we impose:
\begin{align}
\sum w_i &= 1 \quad \text{(constant reproduction)}, \label{eq:const} \\
\sum w_i \Delta x_i &= 0 \quad \text{(cancel $u_x$ term)}, \label{eq:linear} \\
u_t \sum w_i \Delta t_i + \frac{1}{2} u_{xx} \sum w_i \Delta x_i^2 &= 0 \quad \text{(cancel first-order error)}. \label{eq:error}
\end{align}

Substituting $u_t = \alpha u_{xx}$ into \eqref{eq:error}:
\[
\alpha \sum w_i \Delta t_i + \frac{1}{2} \sum w_i \Delta x_i^2 = 0.
\]

\subsection*{Evaluating the Sums}
\begin{align*}
\sum w_i \Delta x_i &= -\Delta x \, w_1 + 0 + \Delta x \, w_3 = \Delta x(w_3 - w_1), \\
\sum w_i \Delta t_i &= -\Delta t(w_1 + w_2 + w_3) = -\Delta t, \\
\sum w_i \Delta x_i^2 &= \Delta x^2(w_1 + w_3).
\end{align*}

\subsection*{Solving the System}
From \eqref{eq:linear}: $w_3 = w_1$.

From the PDE constraint:
\[
\alpha \cdot (-\Delta t) + \frac{1}{2} \cdot \Delta x^2(w_1 + w_3) = 0.
\]
Using $w_3 = w_1$:
\[
-\alpha \Delta t + w_1 \Delta x^2 = 0 \quad \Rightarrow \quad w_1 = \frac{\alpha \Delta t}{\Delta x^2}.
\]

From \eqref{eq:const}:
\[
w_2 = 1 - 2w_1 = 1 - \frac{2\alpha \Delta t}{\Delta x^2}.
\]

\subsection*{Solution}
\[
\boxed{
w_1 = w_3 = \frac{\alpha \Delta t}{\Delta x^2}, \quad
w_2 = 1 - \frac{2\alpha \Delta t}{\Delta x^2}.
}
\]

\subsection*{Explicit Scheme}
\begin{align*}
u_j^n &= w_1 u_{j-1}^{n-1} + w_2 u_j^{n-1} + w_3 u_{j+1}^{n-1} \\
&= \frac{\alpha \Delta t}{\Delta x^2} u_{j-1}^{n-1} + \left(1 - \frac{2\alpha \Delta t}{\Delta x^2}\right) u_j^{n-1} + \frac{\alpha \Delta t}{\Delta x^2} u_{j+1}^{n-1}.
\end{align*}

Rearranging:
\begin{equation}
\boxed{
u_j^n = u_j^{n-1} + \frac{\alpha \Delta t}{\Delta x^2} (u_{j+1}^{n-1} - 2u_j^{n-1} + u_{j-1}^{n-1}).
} \label{eq: fDStencil}
\end{equation}


This is the classical forward Euler scheme for $u_t = \alpha u_{xx}$.

\subsection*{Stability}
For non-negative weights:
\[
w_2 \geq 0 \quad \Rightarrow \quad \boxed{\frac{\alpha \Delta t}{\Delta x^2} \leq \frac{1}{2}.}
\]
This is the CFL stability condition.

For a 5-point symmetric stencil with the same PDE constraint, the
method recovers the standard 4th-order FD scheme computationally (see
\cref{sec:fd4_recovery}).

%% ============================================================
\chapter{Reconstruction Environment}
\label{chap:reconstruction}
%% ============================================================
\section{Overview}
\label{sec:recon_overview}
The reconstruction environment is defined as follows: given a growing cloud of evaluated space time points (i.e. neighbours), can it compute the solution at a new query point (i.e. the target point)? It will eventually serve as the step function for the RL agent, but is already fully operational as a standalone solver.
The design follows a space time formulation: both $x$ and $t$ are treated as unstructured dimensions.
Time is not marched in regular layers; instead, points are added to the cloud. Also, it is worth noting that the following notes take the diffusion advection equation as an example for the algorithm and later the numerical results.

\section{Algorithm: Space--Time GFDM Rollout}
\label{alg:algorithm}
\begin{algorithm}
\caption{Space--Time GFDM Rollout (current implementation)}
\begin{algorithmic}[1]
\State \textbf{Input:} PDE parameters $(\alpha,c)$; reference grid $(x_{\mathrm{FD}},t_{\mathrm{FD}})$; stencil size $n$; maximum number of accepted points $N_{\max}$.
\State \textbf{Initialise:} evaluate the initial condition at all points $(x_i,0)$ and insert them into the visited set.
\State Insert boundary-condition points $(0,t_j,0)$ and $(1,t_j,0)$ for all reference times $t_j$.
\State Build an initial "candidate" set from reference-grid points in a local forward box above the initial condition where the .
\While{candidate set non-empty and number of accepted interior points $< N_{\max}$}
    \State Sample randomly a candidate target point from the candidate set (will be replaced with RL later).
    \If{$z^*$ is too close to an already visited point}
        \State reject and continue
        \State\textbf{ might remove this}
    \EndIf
    \State Query previously visited points with $t_i < t^*$ using an expanding search in normalised space time coordinates (this should work for any time later).
    \If{fewer than $n$ admissible neighbours are found (max box size)}
        \State reject and continue
    \EndIf
    \State Form the offsets $\Delta x_i = x_i-x^*$ and $\Delta t_i = t_i-t^*$. 
    \State Assemble the PDE-informed constraint matrix $A$ and right-hand side $b$.
    \State Solve for the minimum-norm weight vector $w$.
    \If{$A$ is ill-conditioned or $w$ fails the weight admissibility tests}
        \State reject and continue
    \EndIf
    \State Compute the predicted value $\hat u^* = \sum_{i=1}^n w_i u_i$. 
    \State Spawn new forward candidate points from a local box around $z^*$.
\EndWhile
\State \textbf{Output:} visited space--time cloud with predicted solution values.
\end{algorithmic}
\end{algorithm}

\begin{remark}
The neighbour search box is \emph{normalised}: the spatial and
temporal extents are scaled separately because $\Delta x$ and
$\Delta t$ have different natural sizes.
The box is expanded geometrically until $n$ neighbours are found
or a maximum scale is exceeded.
\end{remark}

\begin{remark}
The use of FD in the background is a simplification at this stage. It is used for comparison and candidate choosing.
It means the method is not fully mesh-free yet.
\end{remark}
\pagebreak
\section{Constraint Matrix Choices}
During implementation, it was noticed that weight construction from the matrix is not sufficient for a stable rollout.
Thus, we define two additional constraints:
\begin{enumerate}
    \item Geometric conditions
    \item weight conditions 
\end{enumerate}
Thus, the local predictor is determined not only by the choice of rows in $A$, but also by a stencil acceptance policy. Such conditions will help during the implementation of the RL part, as they enable the agent to learn which neighbours to choose. 
\subsection{Geometric admissibility of the stencil}

Let $z^* = (x^*,t^*)$ be a target point, and let
\[
\mathcal N(z^*) = \{z_i = (x_i,t_i)\}_{i=1}^n
\] 
For now, the rollou only takes neighbours from the past time, i.e.
\[
\Delta t_i < 0, \qquad i=1,\dots,n.
\]
The local reconstruction is attempted only if $\mathcal N(z^*)$ satisfies a geometric criterion. It is worth noting that the following constraints make the method less robust and very sensitive to the geometrical structure of the stencil. This is why other methods like Radial Basis Functions-FD and Generalized Moving Least Squares are yet to be discussed and compared to this one. 

Define the spatial and temporal local scales
\[
h_x = \max_{neighbours} |\Delta x_i|, \qquad
h_t = \max_{neighbours} |\Delta t_i|.
\]
The stencil is declared geometrically admissible if:

\paragraph{Two-sided spatial coverage.}
For interior reconstructions, the neighbour set must contain points on both sides of the target in the spatial direction:
\[
\min_i \Delta x_i < 0 < \max_i \Delta x_i.
\] 
This excludes one-sided interior stencils, which behave as extrapolation, and we noticed that those with such geometries perform poorly. Please note that 3 other constraints were tested, related to proximity (which will be later taken care of with RL), the distance between neighbours in the spatial dimension, and the condition number of the Gram Matrix of the set $dx_i$ and $dt_i$. It is worth noting that all of them were relaxed and didn't affect the error of our rollout, thus, they were removed. 

\subsection{Weight admissibility and stability constraints}

Given an admissible stencil $\mathcal N(z^*) \in \mathcal A_{\mathrm{geom}}$, the reconstruction weights are computed from a PDE-informed constrained system. For the advection--diffusion equation
\begin{equation}
    u_t + c\,u_x = \alpha u_{xx}, \label{diff_advection}
\end{equation}
the implementation uses the normalised matrix
\[
A(\mathcal N) =
\begin{pmatrix}
1 & 1 & \cdots & 1 \\
\dfrac{\Delta x_1 - c\Delta t_1}{h_x} &
\dfrac{\Delta x_2 - c\Delta t_2}{h_x} &
\cdots &
\dfrac{\Delta x_n - c\Delta t_n}{h_x} \\
\dfrac{\frac12\Delta x_1^2 + \alpha \Delta t_1}{h_x^2} &
\dfrac{\frac12\Delta x_2^2 + \alpha \Delta t_2}{h_x^2} &
\cdots &
\dfrac{\frac12\Delta x_n^2 + \alpha \Delta t_n}{h_x^2} \label{A_matrix}
\end{pmatrix},
\qquad
b=
\begin{pmatrix}
1\\0\\0
\end{pmatrix}.
\]
Thus the weight vector $w\in\mathbb R^n$ is defined by the least-squares/minimum-norm solution
\[
w = \arg\min_{\tilde w \in \mathbb R^n}\|\tilde w\|_2
\quad
\text{subject to}
\quad
A(\mathcal N)\tilde w \approx b,
\]
which numerically is computed through
\[
w = \mathrm{lstsq}(A,b).
\]
Since $A\in\mathbb R^{3\times n}$ and the implementation uses $n=5$, this is an underdetermined system with a minimum-norm solution.

The resulting weight vector is accepted only if it satisfies the following stability constraints.

\paragraph{Partition of unity.}
The constant-reproduction constraint is rechecked numerically by requiring
\[
\left|\sum_{i=1}^n w_i - 1\right| \le 10^{-6}.
\]
This guarantees that the reconstruction remains a consistent value predictor after the numerical solve. Removing it didn't affect the number of points taken; however, it was kept as a sanity check, and it is computationally cheap

\paragraph{No dominant coefficient.}
To prevent the local predictor from collapsing onto a single neighbour, the largest weight magnitude is required to satisfy
\[
\max_{1\le i\le n}|w_i|
\;\le\;
\frac{\gamma_w}{n},
\]
where in the implementation
\[
\gamma_w = 3.
\]
For $n=5$, this gives the explicit bound
\[
\max_i |w_i| \le 0.6.
\]
This criterion suppresses highly localised weight patterns, which in practice are associated with unstable extrapolatory behaviour.

\paragraph{Lower bound on negative weights.}
The weight vector is further required to satisfy
\[
\min_{1\le i\le n} w_i \ge w_{\min},
\qquad
w_{\min}=-1.
\]
This does not enforce positivity, but excludes excessively negative coefficients, which tend to produce oscillatory cancellation and error amplification.

 
for which the computed solution $w$ survive the mentioned constraints. The predicted value is then
\[
\hat u(z^*) = \sum_{i=1}^n w_i u(z_i).
\]
 

\chapter{Numerical Experiments}
\label{chap:experiments}
\section{Exp 1: Stencil Accuracy vs.\ \# of Neighbours}
\label{sec:exp1}

\subsection{Setup}

We fix $m=3$ Taylor terms (constant, $\Delta x$, $\Delta t$) and
vary the number of neighbours $n \in \{2, 3, 6, 10\}$.
For each $n$, we draw neighbours uniformly at random in a disk of
radius $r = 0.1$ around a fixed target, compute weights via the
pseudoinverse, and measure the average prediction error against
the true solution $u(x,t) = \sin(\pi x)\cos(\pi t)$ over 100
random targets.

\subsection{Results}

\begin{figure}[h]
    \centering
    \includegraphics[width=0.55\textwidth]{Figures/error_vs_neighbours.png}

    \caption{Average prediction error as a function of the number
    of neighbours $n$, with $m=3$ Taylor constraints.
    Error drops sharply from $n=2$ to $n=3$, then saturates near
    $2\times10^{-3}$ for $n \ge 6$, confirming that the
    underdetermined regime ($n > m$) is optimal.}
    \label{fig:error_vs_nb}
\end{figure}

\subsection{Discussion}

Following the empirical observation of \cite{benito2001gfdm}, a ratio of $n=2m$ neighbours to Taylor constraints is generally sufficient; beyond this point the minimum-norm weights stabilise and additional neighbours bring negligible accuracy gain, a trend confirmed by Figure 4.1

\section{Experiment 2: Error Bound Verification}
\label{sec:exp2}

We verify the inequality (\ref{eq: error_bound}) numerically.
For $m=3$, $n=5$ and the test function $u(x,t) = \sin(\pi
x)\cos(\pi t)$ (so $C = \pi^2 |\sin(\pi x^*)\cos(\pi t^*)|$), it is worth mentioning that no PDE constraints were added to the matrix $A$:
\begin{itemize}
    \item \textbf{Fixed target, shrinking $h$}: both the true error
    and the bound scale as $O(h^2)$, confirming the convergence rate.
    \item \textbf{Random targets at fixed $r=0.1$}: over 80 random
    target points, the true error lies below the bound in all cases.
\end{itemize}
The bound is tight to within a factor of $\sim 10$ for typical
geometries.
\begin{figure}[htbp]
    \centering
    \includegraphics[width=\textwidth]{Figures/error_prop.png}
    \caption{Error bound verification for $m=3$, $n=5$,
    $u(x,t)=\sin(\pi x)\cos(\pi t)$.\\
    \textbf{Left:} example point cloud showing 5 neighbours within
    disk of radius $r=0.12$ around a target point, with effective
    radius $h=0.114$.\\
    \textbf{Centre:} true error and bound vs.\ $h$ for a fixed target
    with shrinking $r$, comparing fixed geometry (same directions
    scaled) and random geometry (fresh sample at each $r$); both
    error curves are parabolic, confirming $O(h^2)$ convergence
    regardless of geometry.\\
    \textbf{Right:} scatter of true error vs.\ bound over 80 random
    target points at fixed $r=0.1$; all points lie below the diagonal,
    confirming the bound holds in every case.
    The bound is tight to within a factor of ${\sim}10$ for typical
    geometries.}
    \label{fig:exp2_error_bound}
\end{figure}
\section{Experiment 3: Recovery of Standard FD Stencils}
\label{sec:exp3}

\subsection{2nd-order (FTCS) stencil}
\label{sec:ftcs_recovery}

We try to verify the theoretical results obtained in (\ref{eq: fDStencil}) with $n=3$ neighbours
on a regular grid, the recovered weights match the classical FTCS
weights to machine precision.
Weight comparison is shown in \cref{fig:weights_2nd}.

\begin{figure}[h]
    \centering
    \includegraphics[width=0.7\textwidth]{Figures/2nd_order_method.png}
    \caption{Left: weight comparison (our method vs.\ FD) for the
    3-point stencil --- values agree exactly.
    Centre: solution snapshot at $t=0.202$.
    Right: mean error over time; our method vs.\ FD vs.\ true solution.
    The error between our method and FD is at the level of
    $10^{-14}$, confirming exact recovery.}
    \label{fig:weights_2nd}
\end{figure}

\subsection{4th-order stencil}
\label{sec:fd4_recovery}

With $n=5$ neighbours,
the recovered weights match the standard 4th-order FD coefficients
to machine precision (\cref{fig:weights_4th}).
At boundaries, the stencil is skewed: all five neighbours are taken
on one side, maintaining 4th-order accuracy without ghost points.

\begin{figure}[h!]
    \centering
    \includegraphics[width=0.7\textwidth]{Figures/4th_order_method.png}
    \caption{Same layout as \cref{fig:weights_2nd} for the 5-point
    stencil.
    The centre weight is essentially zero (the stencil has no
    self-coupling), consistent with the standard 4th-order FD
    formula.
    Error between our method and FD is at the level of $10^{-13}$.}
    \label{fig:weights_4th}
\end{figure}
\pagebreak


\section{Experiment 4: Full Space--Time Rollout}
\label{sec:exp4}

\subsection{Setup}


We run \cref{alg:algorithm} on the advection--diffusion equation
\eqref{diff_advection} using the constraint matrix \eqref{A_matrix}.
A fine FD grid serves both as the reference solution and as the pool
of candidate evaluation locations. The run produces 30{,}000 solution points, evaluated against the true analytical solution.

\subsection{Results}

\begin{table}[h]
    \centering
    \begin{tabular}{ll}
        \toprule
        Metric & Value \\
        \midrule
        Points computed & 30{,}000 \\
        Mean error (vs.\ true) & $3.15 \times 10^{-4}$ \\
        Max error (vs.\ true) & $2.79 \times 10^{-3}$ \\
        Maximum time reached & $t_{\max} = 0.223$ \\
        Total attempts & 133{,}722 \\

        \bottomrule
    \end{tabular}
    \caption{Full rollout performance on the advection--diffusion equation,
    evaluated against the analytical solution.}
    \label{tab:rollout_results}
\end{table}


\begin{figure}[h]
    \centering
    
    \includegraphics[width=0.8\textwidth]{Figures/Figure 3.png}
    
    \caption{Left: spatial temporal cloud of 30{,}000 visited points coloured
    by pointwise error $|u - u_{\mathrm{true}}|$. The error is uniformly
    small across the domain.
    Right: solution slice at $t = 0.0183$.}
    \label{fig:rollout_visited}
\end{figure}

\begin{figure}[h]
    \centering
    \includegraphics[width=0.8\textwidth]{Figures/Figure 4.png}
    
    \caption{Zoomed view of a representative stencil at point
    $(0.343, 0.1174)$. The five neighbours (blue squares) are
    scattered below the target (red circle),
    with the search box shown in pink.}
    \label{fig:rollout_stencil}
\end{figure}


Out of 133{,}722 total attempts, 30{,}000 points were successfully
computed. The rejection breakdown is as follows:

\begin{table}[h!]
    \centering
    \begin{tabular}{lrr}
        \toprule
        Rejection reason & Count & Fraction \\
        \midrule
        No neighbours found & 10 & 0.0\% \\
        Geometry (ill-conditioned $\mathbf{A}$) & 17{,}776 & 13.3\% \\
        Condition number threshold & 46{,}876 & 35.1\% \\
        Weight norm threshold & 39{,}060 & 29.2\% \\
        \bottomrule
    \end{tabular}
    \caption{Rejection breakdown for the full rollout.
    The dominant failure modes are poor conditioning of the local
    system (35\%) and large weight norms (29\%), both of which are
    direct manifestations of the error bound \eqref{sec:error_bound}:
    large $\mathrm{cond}(\mathbf{A})$ and large $\|w\|_1$ indicate
    stencils that would amplify the reconstruction error.}
    \label{tab:rejection}
\end{table}
\begin{figure}[h]
    \centering
    \includegraphics[width=\textwidth]{Figures/Figure 5.png}
    \caption{Left two panels: total accumulated error and running average
    error as a function of number of points evaluated, comparing the
    present method (GFDM, blue) and FD (orange) against the analytical
    solution. Right two panels: scatter of per-point error
    vs.\ condition number of the local $\mathbf{A}$ matrix and vs.\
    $\ell_1$ norm of the weights. In both cases the present method
    (blue) clusters systematically below FD (orange), consistent with
    the error bound \eqref{eq: error_bound}.}
    \label{fig:rollout_errors}
\end{figure}
\subsection{Discussion}

The dominant failure modes --- conditioning and weight norm rejections
--- are not arbitrary numerical thresholds. They are direct consequences
of the error bound \eqref{eq:error_bound}: the bound predicts that
large $\|w\|_1$ amplifies the remainder $Ch^2$, so rejecting high
weight-norm stencils is a principled quality control step, not a
heuristic filter. The fact that 77{,}936 out of 133{,}722 attempts
(58\%) were rejected on these grounds points to a clear direction for
improvement: the current random candidate selection wastes a large
fraction of attempts on geometrically poor stencils. This is precisely
the gap the RL agent is intended to fill --- by learning to select
candidates that are likely to produce well-conditioned stencils with
small weight norms, the agent can dramatically reduce the rejection
rate and improve wall time. The 1499 s wall time for 30{,}000 points is a known bottleneck and will be addressed by replacing the
expanding-box neighbour search with a $k$-d tree spatial index.\chapter{Reinforcement Learning Formulation}

\section{Motivation for RL-Based Point Placement}


So far, what we have tested is the GFDM system that works on random choice of particle movement; what we need to now is to form the RL environment to see whether the reinforcement learning agent will be able to discover the integration schemes that emerge from the physics without feeding the agent any details about the physics. 

First thing to do was to set up the environment (Gymnasium, and PDE) to be able to choose how the policy works and whole setup will be like.

First thing we did was to build the initial


The reconstruction environment developed in Chapter~\ref{chap:reconstruction} is fully
operational as a standalone solver, but the candidate point selection it currently employs
is purely random: at each step, a target point is drawn uniformly from a pre-defined
finite-difference grid and accepted only if it passes the geometric and weight admissibility
tests of Section~\ref{sec:admissibility}. This strategy is demonstrably wasteful, as shown
in Table~\ref{tab:rejection}, but wasteful is not the core problem --- \emph{uninformed}
is.

The deeper issue is that random selection is blind to the causal structure of the PDE.
For any target point $z^* = (x^*, t^*)$, the accuracy of the GFDM prediction depends
critically on the quality of the support structure built below it in space--time. A
collection of evaluation points that respects the domain of dependence of the PDE will
produce well-conditioned stencils and accurate predictions; a collection that ignores it
will not, regardless of how many points are evaluated.

This motivates the central question of the RL component: \emph{given a target point
$z^*$, can an agent learn, purely from the terminal prediction error, an optimal strategy
for placing intermediate evaluation points in space--time?}

The answer should be physically interpretable. For pure diffusion, the domain of
dependence of $(x^*, t^*)$ is a parabola spreading backward in time, and the optimal
support structure should resemble a tree that fans out spatially as it recedes in $t$. For
pure advection, information travels along a single characteristic $x - ct = \mathrm{const}$,
and the optimal path should collapse to a line. For advection--diffusion, the optimal
structure interpolates between these two extremes, with the fan width governed by the
P\'{e}clet number $\mathrm{Pe} = c / (2\alpha)$. The RL agent is not told any of this; it
must discover the structure from the reward signal alone. Whether it does so --- and
whether the discovered structure matches the theoretical domain of dependence --- is the
primary scientific question.

\section{Problem Setup}

We fix a target point $z^* = (x^*, t^*)$ on the reference grid. A set of $N$ walkers
$\{z_i^{(k)}\}_{i=1}^N$ is initialised at the initial condition:
\begin{equation}
    z_i^{(0)} = (x_i, 0), \quad i = 1, \ldots, N,
\end{equation}
where $\{x_i\}$ are $N$ uniformly spaced nodes on $[0, 1]$. Each walker carries a
solution value $u_i^{(k)}$ computed by the GFDM reconstruction at its current position.

Separately, the algorithm maintains a growing visited list
\begin{equation}
    \mathcal{V}^{(k)} = \{(z_j, u_j)\}
\end{equation}
of all space--time points evaluated so far, including the initial condition and boundary
conditions. The visited list is the data structure used by the GFDM neighbour search and
is entirely decoupled from the RL state: it is not observed by the agent, and its size
grows unboundedly without affecting the state dimension.

An episode terminates when any walker lands within a prescribed tolerance $\varepsilon$
of $z^*$, at which point the GFDM prediction at $z^*$ is computed and the terminal
reward is issued.

\section{MDP Formulation}

We cast the point-placement problem as a finite-horizon Markov Decision Process
$(\mathcal{S}, \mathcal{A}, \mathcal{R}, \mathcal{T}, \gamma)$.

\subsection{State Space}

The state at step $k$ is a fixed-size vector encoding the current position and solution
value of each walker, together with the global target:
\begin{equation}
    s^{(k)} = \bigl[x_1^{(k)},\, t_1^{(k)},\, u_1^{(k)},\;
                    \ldots,\;
                    x_N^{(k)},\, t_N^{(k)},\, u_N^{(k)},\;
                    x^*,\, t^*\bigr]
              \in \mathbb{R}^{3N + 2}.
    \label{eq:state}
\end{equation}

The target coordinates $(x^*, t^*)$ are appended once as a global context, not repeated
per walker, giving state dimension $3N + 2$. This is fixed throughout the episode
regardless of how many points have been added to the visited list. Including $z^*$ in the
state is essential for interpretability: the agent learns a policy \emph{conditioned on
the target}, so after training one can vary $z^*$ systematically and observe how the
learned trajectory structure changes --- collapsing to a characteristic line for
advection-dominated targets, or fanning into a tree for diffusion-dominated ones.

The walker state is updated only when a GFDM solve is accepted. If the solve is
rejected, the walker remains at its previous position and $s^{(k+1)} = s^{(k)}$ for that
walker.

\subsection{Action Space}

At each step $k$, a shared policy $\pi_\theta$ is applied independently to each walker.
For walker $i$, the policy observes the local state tuple $(x_i^{(k)}, t_i^{(k)},
u_i^{(k)}, x^*, t^*)$ and outputs a discrete displacement:
\begin{equation}
    ({\delta x}_i,\, {\delta t}_i) \in
    \{-K_x, \ldots, K_x\} \cdot \Delta x_{\mathrm{FD}}
    \;\times\;
    \{0, 1, \ldots, K_t\} \cdot \Delta t_{\mathrm{FD}},
    \label{eq:action}
\end{equation}
where $K_x = 4$ and $K_t = 4$ in the current implementation, giving $|\mathcal{A}| =
(2K_x + 1)(K_t + 1) = 45$ possible actions per walker. Time displacements are
restricted to $\delta t_i \geq 0$ to enforce causality.

The proposed new position for walker $i$ is:
\begin{equation}
    z_i^{\mathrm{new}} = \bigl(x_i^{(k)} + \delta x_i,\; t_i^{(k)} + \delta t_i\bigr).
\end{equation}

Before the GFDM solve is attempted, $z_i^{\mathrm{new}}$ is checked against the visited
list $\mathcal{V}^{(k)}$. If a point within tolerance $\varepsilon_{\mathrm{col}}$
already exists in $\mathcal{V}^{(k)}$, the move is treated as a collision and rejected
without attempting the solve; the walker state is unchanged. This prevents duplicate
entries in the visited list and wasted evaluations. A displacement of $(\delta x_i, \delta
t_i) = (0, 0)$ is valid and results in an immediate collision rejection, effectively
allowing a walker to stay in place for one step.

The policy is shared across all walkers. Since it is applied independently to each
walker's local state tuple, the number of policy parameters does not grow with $N$, and
the per-walker forward passes are parallelisable.

\subsection{Reward Function}
 
The reward is sparse: no signal is issued during the episode, and a single terminal reward
is given when a walker reaches $z^*$:
\begin{equation}
    r^{(k)} = 0 \quad \text{for all } k < K_{\mathrm{term}},
    \label{eq:dense_reward}
\end{equation}
\begin{equation}
    r^{(K_{\mathrm{term}})} = -\bigl|\hat{u}(z^*) - u_{\mathrm{true}}(z^*)\bigr|
                               - \lambda \, K_{\mathrm{term}},
    \label{eq:terminal_reward}
\end{equation}
where $\hat{u}(z^*)$ is the GFDM prediction at $z^*$ using the visited list at
termination, $u_{\mathrm{true}}(z^*)$ is the analytical solution, $K_{\mathrm{term}}$
is the total number of steps taken to reach $z^*$, and $\lambda > 0$ is a
hyperparameter controlling the accuracy--efficiency tradeoff.
 
The step count $K_{\mathrm{term}}$ includes both accepted and rejected moves. This is
intentional: a rejection wastes a step without advancing the visited list, so the agent
is penalised for attempting geometrically poor locations --- not because rejection is
explicitly discouraged, but because it consumes time budget without contributing to
the support structure around $z^*$.
 
The tradeoff governed by $\lambda$ is the central tension of the problem. A large
$\lambda$ produces a fast but potentially inaccurate agent that rushes toward $z^*$
with minimal intermediate support (the straight-line failure mode). A small $\lambda$
produces a thorough agent that builds a dense support structure at the cost of more
evaluations (the slow dense-coverage mode). The optimal policy at intermediate $\lambda$
is the physically motivated strategy: build just enough support structure to achieve
accurate predictions efficiently. Sweeping $\lambda$ and observing how the learned
trajectory structure changes --- from characteristic-following at high $\lambda$ to
 
\emph{how} to build the support structure --- no stencil quality metrics, no distance
penalties, no shaping terms beyond the step cost. Any trajectory structure that emerges
is therefore a genuine discovery: the agent found it because it leads to accurate terminal
predictions efficiently, not because the reward explicitly encoded it.
 
In deployment, $u_{\mathrm{true}}$ is unavailable. A PDE residual estimate computed
from the local GFDM stencil at $z^*$ will replace it; the design of this estimator is left
to future work.

 
\section{Expected Emergent Structure and Validation}

The primary validation is interpretive, not just numerical. After training, the learned
policy is probed by fixing $z^*$ at various locations and visualising the walker
trajectories in the $x$--$t$ plane. Specifically:

\begin{itemize}
    \item \textbf{Advection-dominated} ($\mathrm{Pe} \gg 1$): trajectories should
          collapse toward the characteristic $x - ct = x^* - ct^*$, with little spatial
          fanning.
    \item \textbf{Diffusion-dominated} ($\mathrm{Pe} \ll 1$): trajectories should fan
          out into a tree structure whose width grows with $\sqrt{\alpha(t^* - t)}$,
          matching the parabolic domain of dependence.
    \item \textbf{Intermediate regime}: fan width should scale with $\mathrm{Pe}^{-1}$,
          interpolating between the two extremes.
\end{itemize}

If the learned trajectories match the theoretical domain of dependence, this constitutes
evidence that the agent has implicitly learned the causal structure of the PDE from
terminal accuracy alone --- a result that goes beyond engineering utility and speaks to
the representational capacity of learned solvers.
 
\chapter{Conclusions and Next Steps  \textcolor{r ed}{TO BE REVIEWED}}
\label{chap:conclusions}
%% ============================================================

\section{Summary of Completed Work}
\label{sec:summary}

The following components are complete and validated:

\begin{enumerate}
    \item \textbf{Taylor-constrained local predictor.}
    Derived from first principles; error bound proved and verified
    numerically.

    \item \textbf{Recovery of classical FD stencils.}
    The 2nd-order (FTCS) and 4th-order FD schemes are exactly
    recovered as special cases of the framework, providing a
    strong consistency check.

    \item \textbf{Full space--time rollout.}
    The algorithm produces 25{,}000 solution points for the
    advection--diffusion equation with mean error
    $2.34 \times 10^{-3}$ against a fine FD reference.

    \item \textbf{Failure mode characterisation.}
    A convergence study at three resolutions identifies the
    dominant failure modes (neighbour scarcity, ill-conditioning,
    weight norm growth) and their dependence on point density.
\end{enumerate}

\section{Immediate Next Steps}
\label{sec:next_steps}

\begin{enumerate}
    \item \textbf{Speed up neighbour search.}
    Replace the expanding-box search with a $k$-d tree or
    ball-tree spatial index to reduce wall time by at least an
    order of magnitude.
 

    \item \textbf{RL environment interface.}
  s

    \item \textbf{UQ formalisation.}
    Extend the preliminary UQ study to quantify uncertainty in the
    solution cloud and use it as an exploration signal for the
    RL agent.

    \item \textbf{RL agent design and training.}
    Begin with a simple policy (greedy on local error estimate)
    as a baseline, then move to a learned policy (e.g.\ DQN or
    PPO on the discretised action space).
\end{enumerate}
\bibliographystyle{plain}
\bibliography{reference}
\end{document}